\documentclass{course_sheet}
\usepackage{amsmath,amssymb}
\usepackage{longtable}


\usepackage{xcolor}
\setlength{\parindent}{0pt}
\setlength{\emergencystretch}{2em}
\begin{document}

\title{Exercise: Adaptive Moment Estimation (Adam)}
\author{}
\date{}
\maketitle
In this exercise, you will apply the \textbf{\textit{Adaptive Moment Estimation (Adam)}} optimization algorithm to a simple concrete problem. You will calculate each step of Adam by hand and observe how the parameter changes over several iterations.

\subsection*{Learning objectives}
By the end of this exercise, you should be able to:

\begin{itemize}
\item
Explain the role of the first and second moments in Adam.

\item
Compute the moving averages of gradients and squared gradients.

\item
Apply bias correction.

\item
Compute an Adam parameter update.

\item
Explain why Adam's update is different from ordinary gradient descent.

\end{itemize}
\clearpage
\section*{Part I - The problem}
Suppose we want to minimize the following function:

\[
f(x) = (x-3)^2
\]

The minimum occurs at

\[
x^* = 3.
\]

The derivative is

\[
\frac{df}{dx} = 2(x-3).
\]

We will use \textbf{\textit{Adam}} to find the minimum.

Assume the following Adam hyperparameters:

\[
\alpha = 0.1
\]

\[
\beta_1 = 0.9
\]

\[
\beta_2 = 0.999
\]

\[
\epsilon = 10^{-8}.
\]

We start with

\[
x_0 = 0.
\]

Initially, Adam's first- and second-moment estimates are both zero:

\[
m_0 = 0,
\qquad
v_0 = 0.
\]

Recall the Adam update rules:

\subsubsection*{Step 1: Compute the gradient}
\[
g_t = \nabla f(x_{t-1})
\]

\subsubsection*{Step 2: Update the first moment}
\[
m_t = \beta_1m_{t-1} + (1-\beta_1)g_t
\]

\subsubsection*{Step 3: Update the second moment}
\[
v_t = \beta_2v_{t-1} + (1-\beta_2)g_t^2
\]

\subsubsection*{Step 4: Correct the bias}
\[
\hat m_t = \frac{m_t}{1-\beta_1^t}
\]

\[
\hat v_t = \frac{v_t}{1-\beta_2^t}
\]

\subsubsection*{Step 5: Update the parameter}
\[
x_t =
x_{t-1}
-
\alpha
\frac{\hat m_t}
{\sqrt{\hat v_t}+\epsilon}.
\]

\clearpage
\section*{Part II - Iteration 1 - Guided calculation}
We begin with

\[
x_0=0.
\]

\subsubsection*{Exercise 2.1 - Compute the gradient}
Calculate

\[
g_1 = 2(x_0-3).
\]

\textbf{\textit{Your answer:}}

\[
g_1 = {\phantom{000}}
\]

\subsubsection*{Exercise 2.2 - Compute the first moment}
Using

\[
m_1 = 0.9m_0 + 0.1g_1,
\]

calculate $m_1$.

\textbf{\textit{Your answer:}}

\[
m_1 = {\phantom{000}}
\]

\subsubsection*{Exercise 2.3 - Compute the second moment}
Using

\[
v_1 = 0.999v_0 + 0.001g_1^2,
\]

calculate $v_1$.

\textbf{\textit{Your answer:}}

\[
v_1 = {\phantom{000}}
\]

\subsubsection*{Exercise 2.4 - Apply bias correction}
Because this is the first iteration,

\[
\hat m_1 =
\frac{m_1}{1-0.9^1}.
\]

Calculate $\hat m_1$.

Then calculate

\[
\hat v_1 =
\frac{v_1}{1-0.999^1}.
\]

\textbf{\textit{Your answers:}}

\[
\hat m_1 = {\phantom{000}}
\]

\[
\hat v_1 = {\phantom{000}}
\]

\subsubsection*{Exercise 2.5 - Update $x$}
Finally, calculate

\[
x_1 =
x_0 -
0.1
\frac{\hat m_1}
{\sqrt{\hat v_1}+10^{-8}}.
\]

\textbf{\textit{Your answer:}}

\[
x_1 \approx {\phantom{000}}
\]

\clearpage
\section*{Part III - Iteration 2 - Work it out yourself}
You should now have a value for $x_1$.

Repeat the same five steps.

\subsubsection*{Exercise 3.1 - Gradient}
\[
g_2 = 2(x_1-3)
\]

\[
g_2 = {\phantom{000}}
\]

\subsubsection*{Exercise 3.2 - First moment}
\[
m_2 = 0.9m_1+0.1g_2
\]

\[
m_2 = {\phantom{000}}
\]

\subsubsection*{Exercise 3.3 - Second moment}
\[
v_2 = 0.999v_1+0.001g_2^2
\]

\[
v_2 = {\phantom{000}}
\]

\subsubsection*{Exercise 3.4 - Bias correction}
\[
\hat m_2 =
\frac{m_2}{1-0.9^2}
\]

\[
\hat v_2 =
\frac{v_2}{1-0.999^2}
\]

Calculate:

\[
\hat m_2 = {\phantom{000}}
\]

\[
\hat v_2 = {\phantom{000}}
\]

\subsubsection*{Exercise 3.5 - Parameter update}
\[
x_2 =
x_1 -
0.1
\frac{\hat m_2}
{\sqrt{\hat v_2}+10^{-8}}
\]

Therefore,

\[
x_2 = {\phantom{000}}
\]

\clearpage
\section*{Part IV - Iteration 3}
\subsubsection*{Exercise 4.1}
Now perform the calculation for a third iteration without the intermediate hints.

Complete the following table.

\begingroup
\small
\setlength{\tabcolsep}{4pt}
\begin{longtable}{l|l}
Quantity & Iteration 3 \\[1.5pt]
\hline
\endfirsthead
Quantity & Iteration 3 \\[1.5pt]
\hline
\endhead
\rule{0pt}{2.4ex}$x_2$ &  $\phantom{-0.00}$ \\[3pt]
\rule{0pt}{2.4ex}$g_3$ &  $\phantom{-0.00}$ \\[3pt]
\rule{0pt}{2.4ex}$m_3$ &  $\phantom{-0.00}$ \\[3pt]
\rule{0pt}{2.4ex}$v_3$ &  $\phantom{-0.00}$ \\[3pt]
\rule{0pt}{2.4ex}$\hat m_3$ &  $\phantom{-0.00}$ \\[3pt]
\rule{0pt}{2.4ex}$\hat v_3$ &  $\phantom{-0.00}$ \\[3pt]
\rule{0pt}{2.4ex}$x_3$ &  $\phantom{-0.00}$ \\[3pt]
\end{longtable}
\endgroup
Use:

\[
g_3 = 2(x_2-3)
\]

\[
m_3 = 0.9m_2+0.1g_3
\]

\[
v_3 = 0.999v_2+0.001g_3^2
\]

\[
\hat m_3=\frac{m_3}{1-0.9^3}
\]

\[
\hat v_3=\frac{v_3}{1-0.999^3}
\]

and

\[
x_3 =
x_2 -
0.1
\frac{\hat m_3}
{\sqrt{\hat v_3}+10^{-8}}.
\]

\clearpage
\section*{Part V - Short conceptual questions}
Answer the following conceptual questions.

\subsubsection*{Exercise 5.1}
What does $m_t$ represent?

Choose the best description:

\begin{itemize}
\item
[ ] The average of the parameters seen so far.

\item
[ ] An exponentially weighted moving average of the gradients.

\item
[ ] An exponentially weighted moving average of the squared parameters.

\item
[ ] The learning rate.

\end{itemize}
Explain your answer in one or two sentences.

\subsubsection*{Exercise 5.2}
What does $v_t$ represent?

\begin{itemize}
\item
[ ] An exponentially weighted moving average of the gradients.

\item
[ ] An exponentially weighted moving average of the squared gradients.

\item
[ ] The current value of the objective function.

\item
[ ] The accumulated parameter updates.

\end{itemize}
\subsubsection*{Exercise 5.3}
Why does Adam use $g_t^2$ when calculating $v_t$?

Consider what happens if a gradient is large in magnitude. How would this affect $v_t$ and, consequently, the size of the parameter update?

\subsubsection*{Exercise 5.4}
Why are the bias-correction terms

\[
1-\beta_1^t
\]

and

\[
1-\beta_2^t
\]

necessary?

\textbf{\textit{Hint:}} Think about the fact that

\[
m_0=v_0=0.
\]


\end{document}
